% Self-weight column (Tutorial 1): background mesh and GIMP domains, in the
% style of GIMP_example2.png. The column is 0.05 x 0.05 x 0.8 m: 16 elements of
% size h = 0.05 m, each holding a 2 x 2 x 2 grid of GIMPs. Only the bottom two
% elements and the top element are drawn; the dots between them stand for the
% 13 elements left out.
% Build: pdflatex column_setup.tex
%        pdftocairo -png -transp -r 300 -singlefile column_setup.pdf column_setup
\documentclass[border=4pt]{standalone}
\usepackage{tikz}
\usetikzlibrary{arrows.meta}

\definecolor{gridgrey}{RGB}{185,185,185}
\definecolor{gapgrey}{RGB}{140,140,140}
\definecolor{activeorange}{RGB}{245,130,32}
\definecolor{mpnavy}{RGB}{0,0,153}
\definecolor{gimpfill}{RGB}{179,179,226}
\definecolor{gimpside}{RGB}{158,158,212}
\definecolor{gimptop}{RGB}{204,204,238}
\definecolor{support}{RGB}{254,238,224}
\definecolor{supportside}{RGB}{250,224,201}
\definecolor{supporttop}{RGB}{255,246,238}
\definecolor{dimpurple}{RGB}{104,36,109}

% heights in units of h
\newcommand{\zb}{2}           % top of the bottom two elements
\newcommand{\zt}{3.2}         % bottom of the top element
\newcommand{\ztop}{4.2}       % top of the column
\newcommand{\gimpside}{0.475} % drawn GIMP side (true side 0.5)

\tikzset{
  mesh/.style={draw=gridgrey, line width=1pt},
  hidden/.style={draw=gridgrey, line width=0.6pt, dash pattern=on 2.5pt off 2pt},
  dimline/.style={draw=dimpurple, line width=1pt},
  dim/.style={dimline, {Latex[length=5pt]}-{Latex[length=5pt]}},
  gimpedge/.style={draw=mpnavy, line width=0.5pt, line join=round},
  activenode/.style={fill=activeorange, draw=black, line width=0.5pt},
}

\newcommand{\activenode}[1]{\path[activenode] #1 circle (3pt);}
\newcommand{\mpoint}[1]{\fill[mpnavy] #1 circle (1.4pt);}
% three dots at (#2,#3) across the gap, standing for the elements left out;
% lowered with depth #3 so that they stay in the gap seen in front of them
\newcommand{\gapdots}[3]{%
  \foreach \dz in {-0.25,0,0.25} {\fill[#1] (#2,#3,{2.6-0.2*#3+\dz}) circle (1.2pt);}}

% x and y in [0,1]
\newcommand{\hiddenmesh}{%
  \draw[hidden] (0,1,0) -- (0,1,\zb);
  \draw[hidden] (0,1,\zt) -- (0,1,\ztop);
  \gapdots{gridgrey}{0}{1}
  \foreach \z in {0,1,\zt} {
    \draw[hidden] (0,0,\z) -- (0,1,\z) -- (1,1,\z);
    \path[activenode, opacity=0.45] (0,1,\z) circle (3pt);}}
\newcommand{\visiblemesh}{%
  \foreach \x/\y in {0/0,1/0,1/1} {
    \draw[mesh] (\x,\y,0) -- (\x,\y,\zb);
    \draw[mesh] (\x,\y,\zt) -- (\x,\y,\ztop);
    \gapdots{gapgrey}{\x}{\y}}
  \foreach \z in {0,1,\zt} {\draw[mesh] (0,0,\z) -- (1,0,\z) -- (1,1,\z);}
  \foreach \z in {\zb,\ztop} {\draw[mesh] (0,0,\z) -- (1,0,\z) -- (1,1,\z) -- (0,1,\z) -- cycle;}
  \foreach \z in {0,1,\zb,\zt,\ztop} {
    \activenode{(0,0,\z)} \activenode{(1,0,\z)} \activenode{(1,1,\z)}}
  \activenode{(0,1,\zb)} \activenode{(0,1,\ztop)}}
\newcommand{\dimensions}{%
  \draw[dimline, {Latex[length=5pt]}-] (-0.4,0,0) -- (-0.4,0,2.1);
  \draw[dimline, -{Latex[length=5pt]}] (-0.4,0,3.1) -- (-0.4,0,\ztop);
  \gapdots{dimpurple}{-0.4}{0}
  \node[left=3pt, font=\sffamily] at (-0.4,0,2.6) {0.8 m};
  \draw[dim] (0,0,-0.35) -- node[below=1pt] {$h$} (1,0,-0.35);}
% the support: the visible faces of the elements drawn
\newcommand{\supportblock}[2]{%
  \fill[support] (0,0,#1) -- (1,0,#1) -- (1,0,#2) -- (0,0,#2) -- cycle;
  \fill[supportside] (1,0,#1) -- (1,1,#1) -- (1,1,#2) -- (1,0,#2) -- cycle;}

% one GIMP domain, lower corner (#1,#2,#3), side #4: the material point at the
% centre, then the three visible faces
\newcommand{\gimp}[4]{%
  \mpoint{(#1+#4/2,#2+#4/2,#3+#4/2)}
  \filldraw[gimpedge, fill=gimpfill, fill opacity=0.55]
    (#1,#2,#3) -- (#1+#4,#2,#3) -- (#1+#4,#2,#3+#4) -- (#1,#2,#3+#4) -- cycle;
  \filldraw[gimpedge, fill=gimpside, fill opacity=0.55]
    (#1+#4,#2,#3) -- (#1+#4,#2+#4,#3) -- (#1+#4,#2+#4,#3+#4) -- (#1+#4,#2,#3+#4) -- cycle;
  \filldraw[gimpedge, fill=gimptop, fill opacity=0.55]
    (#1,#2,#3+#4) -- (#1+#4,#2,#3+#4) -- (#1+#4,#2+#4,#3+#4) -- (#1,#2+#4,#3+#4) -- cycle;}

\begin{document}
% legend glyphs, drawn outside the main picture so they use plain axes
\newsavebox{\legmp}\sbox{\legmp}{\tikz[baseline=-0.6ex]\fill[mpnavy] (0,0) circle (1.4pt);}
\newsavebox{\leggimp}\sbox{\leggimp}{\tikz[baseline=-0.6ex]\filldraw[gimpedge, fill=gimpfill, fill opacity=0.55] (0,-0.15) rectangle (0.3,0.15);}
\newsavebox{\legnode}\sbox{\legnode}{\tikz[baseline=-0.6ex]\path[activenode] (0,0) circle (3pt);}
\newsavebox{\legsupport}\sbox{\legsupport}{\tikz[baseline=-0.6ex]\fill[support] (0,-0.15) rectangle (0.3,0.15);}

\begin{tikzpicture}[x={(1.5cm,0cm)}, y={(0.75cm,0.6cm)}, z={(0cm,1.5cm)}]

  % (a) background mesh, with the elements shaded as the support
  \begin{scope}
    \supportblock{0}{\zb}
    \supportblock{\zt}{\ztop}
    \foreach \z in {\zb,\ztop} {\fill[supporttop] (0,0,\z) -- (1,0,\z) -- (1,1,\z) -- (0,1,\z) -- cycle;}
    \hiddenmesh
    \visiblemesh
    \dimensions
    \node[font=\sffamily] at (0.65,0,-1.0) {(a) background mesh};
  \end{scope}

  % (b) a 2 x 2 x 2 grid of GIMP domains in each element, drawn back to front
  \begin{scope}[xshift=5cm]
    \hiddenmesh
    \pgfmathsetmacro{\gap}{(0.5-\gimpside)/2}
    \foreach \j in {1,0} {
      \foreach \i in {0,1} {
        \foreach \z in {0,0.5,1,1.5,\zt,\zt+0.5} {
          \gimp{0.5*\i+\gap}{0.5*\j+\gap}{\z+\gap}{\gimpside}}}}
    \foreach \x in {0.25,0.75} {\gapdots{mpnavy!45}{\x}{0.75}} % back GIMP columns
    \foreach \x in {0.25,0.75} {\gapdots{mpnavy}{\x}{0.25}}    % front GIMP columns
    \visiblemesh
    \dimensions
    \node[font=\sffamily] at (0.65,0,-1.0) {(b) GIMP domains};
  \end{scope}

  % legend, centred under the two panels
  \node[font=\sffamily, anchor=north] at (3.45cm,-1.95cm) {%
    \usebox{\legmp}~material point\hspace{1.6em}%
    \usebox{\leggimp}~GIMP domain\hspace{1.6em}%
    \usebox{\legnode}~active node\hspace{1.6em}%
    \usebox{\legsupport}~support};

\end{tikzpicture}
\end{document}
